\cotefile{93}
\scanpage{92}
% archive foliation: 91
Interprétant $c_{2,3}$ comme
\[
c_{\unsure{2|3},3}:\Omega\unsure{\longrightarrow|\mathbin{-}}\underline{\omega}=\Omega^1_{X/S}
\]
on voit que c'est aussi, compte tenu de
\[
\Omega\simeq\sigma^*(\underline{\Omega}^1_{P/X})
\]
la dérivée relative à la connexion $c$ de la section $\sigma$. on dit (si $\underline{\omega}$ est loc. libre de \unsure{t.f.}) que $\sigma$ (i.e. $\mathcal{L}\subset\mathcal{H}$) est \uline{nulle part stationnaire} pour $c$ si $c_{2,3}$ est ``partout non nulle''. \deleted{\unsure{Appl.}} \inserted{\unsure{i.e.}} $c$ fait de $\Omega$ un facteur direct local de $\underline{\omega}$.

On suppose pour la suite, qu'il en est ainsi, \uline{et} que $\omega$ est inv., \unsure{i.e. encore} que $c_{2,3}:\Omega\longrightarrow\underline{\omega}$ est un iso. On appelle ces connexions ``\uline{transverses} à $\sigma$ (ou à $\mathcal{L}\subset\mathcal{H}$)). On identifie $\Omega\simeq\underline{\omega}$ par $c_{2,3}$, et de façon précise, on suppose donné un iso
\[
\Omega\xrightarrow[\sim]{i}\underline{\omega}=\Omega^1_{X/S}
\]
et on se borne aux conn. pour lesquelles $c_{2,3}=i$. \deleted{\unsure{Elles forment}} \inserted{\unsure{Ce sont les}} sections d'un torseur sous $\mathcal{L}\otimes\underline{\omega}$ [soit $\underline{\operatorname{Conn}}^{\mathrm{trans}}(\mathcal{H},i)$].

Considérons la restriction d'une connexion $c$ à $\Omega\simeq\underline{\omega}=\operatorname{Fil}_1\mathcal{H}$
\[
c|\underset{\substack{\Vert\\\operatorname{Fil}_1\mathcal{H}}}{\Omega}:\quad
\Omega\longrightarrow\underbrace{\mathcal{L}\otimes\omega}_{\unsure{\Omega\otimes\omega,\omega|\Omega^{\otimes2},\omega}}.
\]
